\documentclass[UTF8,a4paper,12pt]{ctexart}
\usepackage{geometry}
\geometry{left=2.4cm,right=2.4cm,top=2.5cm,bottom=2.5cm}
\usepackage{amsmath,amssymb,bm}
\allowdisplaybreaks[2]
\IfFileExists{/home/lijs/.local/share/fonts/NotoSerifSC-Regular.otf}{%
  \setCJKmainfont[Path=/home/lijs/.local/share/fonts/,BoldFont=NotoSerifSC-Bold.otf]{NotoSerifSC-Regular.otf}%
  \setCJKsansfont[Path=/home/lijs/.local/share/fonts/,BoldFont=NotoSansSC-Bold.otf]{NotoSansSC-Regular.otf}%
}{}
\usepackage{newtxtext}
\usepackage{newtxmath}
\usepackage{xcolor}
\usepackage{fancyhdr}
\pagestyle{fancy}
\fancyhf{}
\fancyfoot[C]{\small\thepage}
\renewcommand{\headrulewidth}{0pt}
\renewcommand{\footrulewidth}{0pt}
\usepackage[hidelinks]{hyperref}
\hypersetup{
  pdftitle  = {第 11 章作业（大学物理 II）},
  pdfauthor = {李谨硕},
  pdfsubject = {胡其图教材 习题 11.2、11.3、11.4、11.6},
}

\ctexset{
  section = {format = {\centering\Large\bfseries}},
  subsection = {format = {\large\bfseries}},
}

\begin{document}

\begin{center}
  {\LARGE\bfseries 第 11 章作业}\\[0.55em]
  {\normalsize 大学物理 II（PHY1252H）}\\[0.95em]
  {\small 李谨硕\quad IEEE 计算机方向\quad 电院 2501\quad 2026 年 9 月 23 日}
\end{center}
\vspace{-0.2em}
{\centering\textcolor{black!30}{\rule{0.9\textwidth}{0.8pt}}\par}
\vspace{0.6em}

\section*{胡其图教材}

\subsection*{11.2}

\textbf{（1）细棒中垂线上距离为 $L$ 处的电场强度}

解：直接使用库仑定律计算。由已知
\[
d\vec{E} = k\frac{dq}{r^3}\vec{r},
\qquad
r = \sqrt{\left(x-\frac{L}{2}\right)^2 + L^2}
\]
我们分别求 $x$、$y$ 分量：
\[
\begin{aligned}
E_x &= k\lambda_0\int_0^L \frac{x\left(\frac{L}{2}-x\right)}
      {\left[\left(x-\frac{L}{2}\right)^2+L^2\right]^{3/2}}\,dx \\
&= -k\lambda_0\int_{-L/2}^{L/2}\frac{x^2}{(x^2+L^2)^{3/2}}\,dx \\
&= -2k\lambda_0\left(\ln\frac{1+\sqrt{5}}{2}-\frac{\sqrt{5}}{5}\right)
\end{aligned}
\]
同理不难求得 $E_y$：
\[
E_y = k\lambda_0 L\int_0^L \frac{x\,dx}
      {\left[\left(x-\frac{L}{2}\right)^2+L^2\right]^{3/2}}
= \frac{k\lambda_0}{\sqrt{5}}
\]
因此得到
\[
\boxed{\vec{E} = k\lambda_0\left[
-2\left(\ln\frac{1+\sqrt{5}}{2}-\frac{\sqrt{5}}{5}\right)\hat{\imath}
+ \frac{1}{\sqrt{5}}\hat{\jmath}\right]}
\]

\textbf{（2）细棒延长线上 $2L$ 处的电场强度}

解：由对称性，$E_y = 0$，只需利用积分计算 $E_x$。取
\[
d\vec{E} = k\frac{dq}{r^3}\vec{r},\qquad r = 2L+x
\]
则
\[
E_x = -k\lambda_0\int_0^L\frac{x}{(x+2L)^2}\,dx
= -k\lambda_0\left(\ln\frac{3}{2}-\frac{1}{3}\right)
\]
所以
\[
\boxed{\vec{E} = -k\lambda_0\left(\ln\frac{3}{2}-\frac{1}{3}\right)\hat{\imath}}
\]

\subsection*{11.3}

解：由对称性 $E_x = 0$，只需计算 $E_y$：
\[
d\vec{E} = -kb\sin\phi\,R\,\frac{\vec{r}}{r^3}\cdot d\phi
\]
\[
d\vec{E}_y = -kb\sin\phi\,R\,\frac{1}{R^2}\sin\phi\cdot d\phi
\]
\[
\begin{aligned}
E_y &= -\frac{kb}{R}\int_0^{\pi}\sin^2\phi\,d\phi \\
&= -\frac{kb}{R}\left[\frac{\phi}{2}-\frac{\sin 2\phi}{4}\right]_0^{\pi} \\
&= -\frac{kb\pi}{2R}
\end{aligned}
\]
所以
\[
\boxed{\vec{E} = -\frac{kb\pi}{2R}\hat{\jmath}}
\]

\subsection*{11.4}

先在圆环所在平面内建立平面直角坐标系，再在空间中建立右手直角坐标系：$xOy$ 平面即圆环所在平面，$z$ 轴垂直于圆环所在平面且过圆心，$x$ 轴指向 $\phi=0$ 的方向，$y$ 轴指向 $\phi=\pi/2$ 的方向。

\textbf{（1）圆环轴线上任意一点处的电场强度}

解：由对称性 $E_y = 0$，$E_z = 0$，只需计算 $E_x$。由
\[
d\vec{E} = k\frac{dq}{r^3}\vec{r},\qquad
r = \sqrt{R^2+z^2},\qquad
dq = \lambda(\phi)R\,d\phi = b\cos\phi\,R\,d\phi
\]
可得
\[
E_x = -\int_0^{2\pi}k\frac{b\cos\phi\,R\,d\phi}{R^2+z^2}
      \cdot\frac{R\cos\phi}{\sqrt{R^2+z^2}}
= -\frac{kbR^2}{(R^2+z^2)^{3/2}}\int_0^{2\pi}\cos^2\phi\,d\phi
= -\frac{kb\pi R^2}{(R^2+z^2)^{3/2}}
\]
所以
\[
\boxed{\vec{E} = -\frac{kb\pi R^2}{(R^2+z^2)^{3/2}}\hat{\imath}}
\]

\textbf{（2）证明远处是电偶极子场，并求偶极矩}

解：当 $z\gg R$ 时
\[
\vec{E}\approx-\frac{kb\pi R^2}{z^3}\hat{\imath}
\]
与偶极子的电场公式对比，可得该偶极子的偶极矩为
\[
\boxed{\vec{p} = b\pi R^2\hat{\imath}}
\]

\subsection*{11.6}

解：由已知
\[
\rho(r) = \rho_0\left(1-\frac{r}{R}\right)
\]
计算空间中任意一点的电场强度，需要分类讨论。

当 $r<R$ 时，由高斯定理
\[
\oint\vec{E}\cdot d\vec{A} = \frac{Q_{\text{in}}}{\epsilon_0}
\]
\[
Q_{\text{in}} = 4\pi\rho_0\int_0^r\left(1-\frac{r'}{R}\right)r'^2\,dr'
= 4\pi\rho_0\left(\frac{r^3}{3}-\frac{r^4}{4R}\right)
\]
\[
E(r) = \frac{\rho_0}{\epsilon_0}\left(\frac{r}{3}-\frac{r^2}{4R}\right)
\]

当 $r>R$ 时，先计算总电荷
\[
Q = 4\pi\rho_0\int_0^R\left(1-\frac{r'}{R}\right)r'^2\,dr'
= 4\pi\rho_0\left(\frac{R^3}{3}-\frac{R^3}{4}\right)
= \frac{\pi R^3}{3}\rho_0
\]
\[
E(r) = \frac{Q}{4\pi\epsilon_0 r^2}
= \frac{\pi R^3\rho_0}{12\pi\epsilon_0 r^2}
= \frac{R^3\rho_0}{12\epsilon_0 r^2}
\]

所以在整个空间内，电场为
\[
\boxed{\vec{E}(r) = \begin{cases}
\dfrac{\rho_0}{\epsilon_0}\left(\dfrac{r}{3}-\dfrac{r^2}{4R}\right)\hat{r}, & r<R \\[8pt]
\dfrac{R^3\rho_0}{12\epsilon_0 r^2}\hat{r}, & r\geq R
\end{cases}}
\]

\end{document}
